\originalchapter{8.323 作业3题解}
本章据 MIT OCW 8.323(2023 春)\href{https://ocw.mit.edu/courses/8-323-relativistic-quantum-field-theory-i-spring-2023/mit8_323_s23_pset03.pdf}{作业3}与\href{https://ocw.mit.edu/courses/8-323-relativistic-quantum-field-theory-i-spring-2023/mit8_323_s23_pset_03sol.pdf}{官方解}中文重编,补充并核查推导;不是官方中文版.继续采用 $(+---)$、$[a_{\mathbf k},a^\dagger_{\mathbf q}]=(2\pi)^3\delta^{(3)}(\mathbf k-\mathbf q)$、$\ket{k}=\sqrt{2\omega_{\mathbf k}}a^\dagger_{\mathbf k}\ket0$,$H$ 已减去真空能.对 $\Lambda=\exp\omega$,用 $U_\Lambda=\exp(-\ii\omega_{\mu\nu}M^{\mu\nu}/2)$;这个符号是原稿 $(-+++)$ 的抬降指标约定转换后的结果,保证 $U_\Lambda\ket{k}=\ket{\Lambda k}$.

\originalsection{原题1}
\begin{exercise}
(a) 由产生湮灭算符对易子及作业1原题2(b),说明洛伦兹变换律应为
\[
U_\Lambda a_{\mathbf k}U_\Lambda^\dagger=\sqrt{\frac{\omega_{\Lambda\mathbf k}}{\omega_{\mathbf k}}}a_{\Lambda\mathbf k},\qquad
U_\Lambda a^\dagger_{\mathbf k}U_\Lambda^\dagger=\sqrt{\frac{\omega_{\Lambda\mathbf k}}{\omega_{\mathbf k}}}a^\dagger_{\Lambda\mathbf k}.
\]
$\Lambda\mathbf k$ 指变换后壳上四动量的空间部分.
(b) 用作业2的 $M^{\mu\nu}$ 计算 $\tfrac12[\omega_{\mu\nu}M^{\mu\nu},a_{\mathbf k}]$ 及创造算符版本,验证无穷小变换与(a)一致;(c) 证明有限变换下 $U_\Lambda\ket0=\ket0$,并在接受(a)的前提下证明 $U_\Lambda\ket{k}=\ket{\Lambda k}$,不要求直接把有限指数完全求和.
\end{exercise}
\begin{solution}
(a) 设标量粒子的变换不混合创造和湮灭算符,系数为实正 $c(k)$(相位约定已固定).不变性要求
\[
c(k)c(q)\delta^{(3)}(\Lambda\mathbf k-\Lambda\mathbf q)=\delta^{(3)}(\mathbf k-\mathbf q).
\]
由壳雅可比为 $\omega_{\Lambda k}/\omega_k$,在对角线处得到 $c(k)^2=\omega_{\Lambda k}/\omega_k$.仅由对易子不能排除额外相位;标量表示的相位约定给出题设形式.
(b) 记 $G=\tfrac12\omega_{\mu\nu}M^{\mu\nu}$,$v_i=\delta k_i=\omega^i{}_\nu k^\nu=\omega_{0i}\omega_k-\omega_{ij}k_j$,$\delta\omega_k=\omega_{0i}k_i$.作业2的表达式和 $[a^\dagger_q,a_k]=-(2\pi)^3\delta^{(3)}(q-k)$ 给出
\[
[G,a_k]=\ii\left(v_i\partial_{k_i}+\frac{\delta\omega_k}{2\omega_k}\right)a_k,\qquad
[G,a_k^\dagger]=\ii\left(v_i\partial_{k_i}+\frac{\delta\omega_k}{2\omega_k}\right)a_k^\dagger.
\]
创造算符版本可将第一式取厄米共轭;因 $[G,a]^\dagger=-[G,a^\dagger]$,两式同号.于是 $U aU^\dagger=a-\ii[G,a]+\order(\omega^2)=a+v\cdot\nabla_k a+(\delta\omega_k/2\omega_k)a$.另一方面展开(a)右端的平方根及动量自变量,恰得相同两项.
(c) 正规序 $M$ 的每项含右侧湮灭算符,故 $M^{\mu\nu}\ket0=0$.指数级数除常数项外全消失.再用(a),
$U\ket{k}=\sqrt{2\omega_k}(Ua_k^\dagger U^\dagger)\ket0=\sqrt{2\omega_{\Lambda k}}a^\dagger_{\Lambda k}\ket0=\ket{\Lambda k}$.
\end{solution}

\originalsection{原题2}
\begin{exercise}
考虑 $\ket{\mathbf r,t}_f=\int\dd^3k\,f(\mathbf k)\ket{k}\e^{-\ii\mathbf k\cdot\mathbf r+\ii\omega_k t}$.
(a) 由 ${}_f\bra{\mathbf r_1,t}\mathbf r_2,t\rangle_f=\delta^{(3)}(\mathbf r_1-\mathbf r_2)$ 确定 $f$;(b) 对(a)所得态施加 $U_\Lambda$,证明它一般不等于 $\ket{\Lambda\mathbf r,\Lambda t}_f$;(c) 求类空分离 $r>|t|$ 时的重叠 $C(\mathbf r_1-\mathbf r_2,t_1-t_2)={}_f\bra{\mathbf r_2,t_2}\mathbf r_1,t_1\rangle_f$.若把其模平方视为位置传播概率,是否局限于前向光锥？(d) 比较 $\ket{x}=\phi(x)\ket0$:它是否完全局域？证明其洛伦兹协变性;(e) 设 $\ket\Psi=\int\dd^3k\,h(\mathbf k)\ket{k}/(2\pi)^3$,$\Psi(x)=\bra0\phi(x)\ket\Psi$,求 $h$,使零时刻波函数为中心 $\mathbf x_0$、宽度 $a$、平均动量 $\mathbf p$ 的高斯波包.
\end{exercise}
\begin{solution}
(a) 因 $\braket{q}{k}=2\omega_k(2\pi)^3\delta^{(3)}(\mathbf k-\mathbf q)$,同刻重叠为
$(2\pi)^3\int\dd^3k\,2\omega_k|f(k)|^2\e^{\ii\mathbf k\cdot(\mathbf r_1-\mathbf r_2)}$.故 $|f|^2=[(2\pi)^6 2\omega_k]^{-1}$.选零相位,$f=1/[(2\pi)^3\sqrt{2\omega_k}]$;正交条件本身不确定相位.
(b) 令 $q=\Lambda k$.用 $\dd^3k=\omega_k\dd^3q/\omega_q$ 及 $k\cdot x=q\cdot\Lambda x$,
\[
U\ket{\mathbf r,t}_f=\int\frac{\dd^3q}{(2\pi)^3\sqrt{2\omega_q}}
\sqrt{\frac{\omega_{\Lambda^{-1}q}}{\omega_q}}\ket q\,\e^{\ii q\cdot\Lambda x}.
\]
平方根在一般 boost 下非恒等于1,故不协变.纯旋转则仍协变.
(c) 由定义 $C(r,t)=\int\dd^3k\,\e^{\ii\omega_k t-\ii\mathbf k\cdot\mathbf r}/(2\pi)^3$.引入
$D(r,t)=\int\dd^3k\,\e^{\ii\omega_k t-\ii\mathbf k\cdot\mathbf r}/[(2\pi)^3 2\omega_k]$,则 $C=-2\ii\partial_tD$.$D$ 为洛伦兹标量.令 $s=\sqrt{r^2-t^2}>0$,在同时的参考系中角积分给出
\[
D(s)=\frac1{4\pi^2s}\int_0^\infty\frac{k\sin(ks)}{\sqrt{k^2+m^2}}\dd k
=\frac{m}{4\pi^2s}K_1(ms).
\]
为复核最后一步,利用 $\int_0^\infty\cos(ks)/\sqrt{k^2+m^2}\dd k=K_0(ms)$(令 $k=m\sinh u$ 后的 Bessel 积分表示),对 $s$ 求导得所需正弦积分.$K_\nu$ 是第二类修正 Bessel 函数,满足 $\dd[K_1(ms)/s]/\dd s=-mK_2(ms)/s$.因此
\[
C(r,t)=-\frac{\ii m^2t}{2\pi^2s^2}K_2(ms)\qquad(r>|t|).
\]
当 $t\ne0$ 时它通常非零,所以这种单粒子位置概率解释会给出光锥外尾部.同刻 $t=0,r>0$ 重叠为零,与(a)一致.场论的因果性由类空场对易子为零保证,不能用这一重叠代替.
(d) 创造部分给出 $\ket x=\int\dd^3k\,\e^{\ii k\cdot x}\ket k/[(2\pi)^3 2\omega_k]$.同刻重叠为 $mK_1(m|\mathbf x-\mathbf y|)/(4\pi^2|\mathbf x-\mathbf y|)$,并非 $\delta$.但壳测度和相位均不变,换元 $q=\Lambda k$ 后 $U\ket x=\ket{\Lambda x}$.
(e) $\bra0\phi(x)\ket k=\e^{-\ii k\cdot x}$,故零时刻 $\Psi$ 是 $h$ 的空间逆傅里叶变换.明确取宽度约定
\[
\Psi(\mathbf x,0)=\frac{1}{\pi^{3/4}a^{3/2}}\exp\left[-\frac{|\mathbf x-\mathbf x_0|^2}{2a^2}+\ii\mathbf p\cdot(\mathbf x-\mathbf x_0)\right].
\]
三次使用 $\int\e^{-y^2/(2a^2)-\ii qy}\dd y=\sqrt{2\pi}a\e^{-a^2q^2/2}$,得到
$h(\mathbf k)=2^{3/2}\pi^{3/4}a^{3/2}\e^{-a^2|\mathbf k-\mathbf p|^2/2-\ii\mathbf k\cdot\mathbf x_0}$.
这是 $\int|\Psi|^2\dd^3x=1$ 的波函数约定;它不自动给出 $\braket{\Psi}{\Psi}=1$,后者还含 $2\omega_k$ 的权重.整体相位不影响中心或动量.
\end{solution}

\originalsection{原题3}
\begin{exercise}
对自由实标量场:(a) 证明 $\bra0\phi(x)\ket0=0$,把 $\sigma^2=\bra0\phi(x)^2\ket0$ 写成一维积分并证明发散;(b) 正规序指将所有 $a$ 移到 $a^\dagger$ 右侧,且不在定义中补入对易子.用 $a,a^\dagger$ 表示 $:\phi^2:$,证明其真空均值为零及 $\phi^2=:\phi^2:+\sigma^2\mathbf1$;(c) 定义空间抹平场
\[
\phi_a(\mathbf x,t)=\frac1{\pi^{3/2}a^3}\int\dd^3y\,\e^{-|\mathbf x-\mathbf y|^2/a^2}\phi(\mathbf y,t),\qquad a>0.
\]
证明均值为零,把其方差 $\sigma_a^2$ 写成一维积分并证明有限;(d) 不必精确积分,求小 $a$、大 $a$ 时 $\sigma_a^2\sim\alpha a^\gamma$ 的 $\alpha,\gamma$.
\end{exercise}
\begin{solution}
(a) $a\ket0=0$、$\bra0a^\dagger=0$ 使均值为零;平方中只有 $a_ka_q^\dagger$ 项留下,故
\[
\sigma^2=\intp\frac1{2\omega_p}=\frac1{4\pi^2}\int_0^\infty\frac{k^2\dd k}{\sqrt{k^2+m^2}}.
\]
大 $k$ 被积函数渐近为 $k$,二次发散.
(b) 记 $u_k(x)=\e^{-\ii k\cdot x}$,则
\[
:\phi^2(x):=\int\frac{\dd^3k\dd^3q}{(2\pi)^6\sqrt{2\omega_k2\omega_q}}
[a_ka_qu_ku_q+a_k^\dagger a_q^\dagger u_k^*u_q^*+2a_k^\dagger a_qu_k^*u_q].
\]
每项至少在真空的一侧被湮灭.原平方与它之差为积分中的 $[a_k,a_q^\dagger]u_ku_q^*$,恰等于(a)的常数乘 $\mathbf1$.未调节时的等式是形式表达;用同一截止定义两边后再取去截止极限.
(c) 归一高斯的傅里叶乘子为 $\e^{-a^2k^2/4}$,故均值仍为零,两场的乘子相乘给出
\[
\sigma_a^2=\frac1{4\pi^2}\int_0^\infty\frac{k^2\e^{-a^2k^2/2}}{\sqrt{k^2+m^2}}\dd k.
\]
零端可积,无穷远由高斯压低,所以 $a>0$ 时有限.
(d) 令 $u=ak$,得 $\sigma_a^2=(4\pi^2a^2)^{-1}\int_0^\infty u^2\e^{-u^2/2}/\sqrt{u^2+(ma)^2}\dd u$.小 $a$ 时被积函数受 $u\e^{-u^2/2}$ 控制,支配收敛给出 $\alpha=1/(4\pi^2),\gamma=-2$.固定 $m>0$、大 $a$ 时乘以 $ma$ 后极限为 $u^2\e^{-u^2/2}$,积分是 $\sqrt{\pi/2}$,故
\[
\sigma_a^2\sim\frac1{4\sqrt2\pi^{3/2}m a^3},\qquad \alpha=\frac1{4\sqrt2\pi^{3/2}m},\quad\gamma=-3.
\]
$m=0$ 时所有 $a$ 都是 $1/(4\pi^2a^2)$,不能套用含 $1/m$ 的大尺度式.抹平尺度增加使绝对真空方差趋零.
\end{solution}

\originalsection{原题4}
\begin{exercise}
允许相互作用的复标量理论有相位对称性,且真空也不变:$U_\alpha\ket0=\ket0$.令
$U_\alpha=\e^{-\ii\alpha Q}$,$Q=\ii\int\dd^3x\,(\pi_{\phi^*}\phi^*-\pi_\phi\phi)$.
(a) 不使用仅适用于自由场的模展开,证明 $U_\alpha\phi U_\alpha^\dagger=\e^{\ii\alpha}\phi$、$U_\alpha\phi^* U_\alpha^\dagger=\e^{-\ii\alpha}\phi^*$;(b) 证明 $\bra0\phi(x)\phi(y)\ket0=0$;(c) 证明一般真空关联函数中,若 $\phi$ 与 $\phi^\dagger$ 的数目不同,则关联函数为零.拉格朗日量对称不自动保证真空对称;本题明确采用后一个假设.
\end{exercise}
\begin{solution}
本题沿用原题给定的 $U_\alpha=\e^{-\ii\alpha Q}$,所以这里的 $Q$ 是作业1原题3按直接诺特参数取荷的整体反号;两者都守恒,但生成算符的指数符号必须一起转换.
(a) 同刻正则关系给出 $[Q,\phi]=-\phi$、$[Q,\phi^*]=\phi^*$,交叉场及动量的对易子为零.BCH 重复对易子形成指数,得所求相位.由于 $Q$ 守恒、$[Q,H]=0$,可将该同刻结果演化至任意时刻.
(b) 在关联函数的每个场两侧插入 $U^\dagger U$,并使用真空不变性,得 $F=\e^{2\ii\alpha}F$ 对所有 $\alpha$ 成立.选 $\alpha=\pi/2$ 即 $F=-F$,所以 $F=0$.
(c) 若有 $N$ 个 $\phi$、$M$ 个 $\phi^\dagger$,同样得到 $F=\e^{\ii(N-M)\alpha}F$.当 $N\ne M$ 可选 $\alpha$ 使相位不为1,故 $F=0$.这一论证既不使用微扰展开,也不要求自由理论;相等数目只是非零的必要条件.
\end{solution}
